\title{Parameter-Efficient Orthogonal Finetuning via Factorization}

\begin{abstract}
The proliferation of large foundation models highlights their significance, yet the substantial computational resources required to train them from the beginning present a major barrier. As a result, the focus has shifted toward strategies that enable effective adaptation of these models for new tasks while mitigating costs. This work revisits the framework of Orthogonal Finetuning (OFT) for the purpose of downstream task adaptation. While OFT exhibits robust generalization capabilities, it remains parameter-intensive due to the substantial size of orthogonal matrices involved. To tackle this limitation, we reinterpret OFT through the lens of information transmission, leading us to articulate key criteria for enhancing parameter efficiency. Drawing inspiration from the efficient information routing achieved by the Cooley-Tukey fast Fourier transform algorithm, we introduce an innovative orthogonal parameterization leveraging butterfly structures. This technique is then integrated with OFT, giving rise to Orthogonal Butterfly~(BOFT), a new method for parameter-efficient finetuning. BOFT serves as a generalization of OFT by incorporating it as a particular instance within its broader orthogonal finetuning scheme. To validate our approach, we perform a comprehensive empirical evaluation, adapting large vision transformers, language models, and text-to-image diffusion models across a range of tasks in the domains of vision and language.
\end{abstract}

\section{Introduction}

Cutting-edge models like ChatGPT~\cite{brown2020language,chen2021evaluating} and Stable Diffusion~\cite{rombach2022high} illustrate the extraordinary generalization power inherent to large-scale foundation models. However, these advances are accompanied by a dramatic rise in model parameters—GPT-3, for example, comprises approximately 175B parameters—which renders training such systems from the ground up increasingly impractical. The advancement of the field, therefore, hinges on efficient mechanisms for repurposing existing foundation models without full retraining. This necessity drives the development of methods that enable fast and efficient adaptation to new tasks. The prevailing approaches for this adaptation can be grouped into three categories: \emph{model finetuning}~\cite{hulora2022,zaken2022bitfit,guo2020parameter,raffel2020exploring,qiu2023controlling,zhang2022adaptive,chavan2023one}, where finetuning is restricted to a subset of parameters in an unchanged model architecture; \emph{adapter tuning}~\cite{rebuffi2017learning,houlsby2019parameter,pfeiffer2020adapterfusion,lin2020exploring,he2021towards}, which introduces new trainable modules into the existing model; and \emph{prompt tuning}~\cite{li2021prefix,lester2021power}, where trainable prefixes are appended to the input sequence. Of these, model finetuning is notable for its conceptual simplicity, effective performance, and the advantage of introducing no additional delay at inference time.

The core idea underlying model finetuning is to maintain a close resemblance between the pretrained model and its finetuned counterpart, according to particular similarity metrics, in order to retain the knowledge acquired during pretraining. Conventional finetuning techniques typically utilize a reduced learning rate as a practical heuristic to restrict divergence between the pretrained and finetuned models. Recognizing the inherently structured properties of weight matrices, alternative, theoretically-motivated strategies seek to safeguard relational characteristics—specifically, the angles between neuron vectors. This consideration has inspired the Orthogonal Finetuning~(OFT)~\cite{qiu2023controlling} methodology, wherein neurons within a layer are collectively updated via the same orthogonal transformation, ensuring that pairwise angular relationships remain invariant during finetuning. While OFT achieves compelling results in terms of generalization and optimization for tasks like finetuning text-to-image diffusion models~\cite{qiu2023controlling}, its parameter count tends to be significant, a direct consequence of the large size of orthogonal matrices. In response, OFT incorporates block-diagonal structures to lower the total parameters. Nevertheless, such parameter savings introduce trade-offs: the resulting orthogonal matrices possess a rigid, block-diagonal sparsity pattern, with each block undergoing independent transformations. This design, though efficient in parameter use, may inadvertently lead to unfavorable inductive biases—such as diminished expressivity and an inability to realize arbitrary linear transformations. Additionally, the challenge of identifying an optimal sparsity structure for the orthogonal matrix remains unsolved.

A central challenge lies in constructing a dense orthogonal matrix in a parameter-efficient way. At first, this appears problematic because a fully dense orthogonal matrix in $d$ dimensions naively necessitates $\mathcal{O}(d^2)$ parameters. Instead of parameterizing the matrix directly, we propose an alternative method: assembling the dense orthogonal matrix by successively multiplying several structured sparse matrices. This strategy leverages the observation that one can decrease the parameter count by exchanging additional computation for memory savings; since the matrix is expressed as a product of sparse factors, each training iteration requires repeated computation of their multiplication. Framing this as a matrix factorization issue, we conceptualize dense orthogonal matrix generation as analogous to an information transmission process. From this vantage point, constructing a dense orthogonal matrix via sparse matrix products can be likened to propagating information over a grid-like graph. This communication-inspired viewpoint provides a flexible foundation for designing diverse sparse matrix factorizations that can tightly constrain the number of learnable parameters yet maintain sufficient expressive power to model fully dense orthogonal matrices.

To maximize parameter efficiency within this information transmission paradigm, we are motivated by the butterfly topology featured in the Cooley-Tukey fast Fourier transform algorithm~\cite{cooley1965algorithm}, renowned for enabling rapid information exchange among nodes~\cite{klasing1994broadcasting}. Notably, the butterfly graph of the Cooley-Tukey scheme yields an efficient sparse matrix factorization known as butterfly factorization. Utilizing this construction, we can synthesize a $d$-dimensional dense matrix from the product of only $\mathcal{O}(\log d)$ sparse matrices, each containing just $\mathcal{O}(d)$ nonzero elements. By ensuring orthogonality of every sparse factor, we establish an orthogonal parameterization with a parameter count reduced to merely $\mathcal{O}(d\log d)$, substantially more compact than the standard $\mathcal{O}(d^2)$. Exploiting this efficient scheme, we introduce Orthogonal Butterfly~(BOFT), a novel method for parameter-efficient finetuning. BOFT generalizes OFT as a specific instance and serves as a broad orthogonal finetuning framework. Both the block-diagonal structure of OFT and the butterfly structure implemented in BOFT share the trait of inducing sparsity, thereby limiting the trainable parameter count in the orthogonal matrices. Comparing this approach to the low-rank modeling of LoRA~\cite{hulora2022}, we note that both low-rank and sparse matrices exemplify two principal categories of structured matrices~\cite{candes2011robust} employed to promote parameter efficiency.

In contrast to OFT's block-diagonal framework, which balances model capacity and structural regularity, BOFT adopts the butterfly structure, offering a more continuous transition between matrices from the full orthogonal group (for expressiveness) and the identity matrix (for regularity). This approach allows us to define a more compact subset within the orthogonal group, tailored for downstream tasks. Since the butterfly structure is foundational in numerous efficient linear transformations—such as the discrete Fourier and cosine transforms—we posit that this structured form of parameterization instills a beneficial inductive bias. This bias aids in enhancing generalization while also mitigating the risk of overfitting. This notion is motivated by the observation that the butterfly structure is inherently equipped to approximate various classic linear transformations, a capability lacking in OFT’s block-diagonal arrangement. Our main contributions are as follows:

\begin{itemize}[leftmargin=*,nosep]
    \item We address parameter efficiency in orthogonal finetuning through a new information transmission framework. This perspective recasts the challenge of constructing dense, parameter-efficient orthogonal matrices as a problem of information propagation on a grid-based graph.
    \item Drawing inspiration from the butterfly configurations used in the Cooley-Tukey algorithm, we introduce Orthogonal Butterfly—a parameter-efficient finetuning technique for orthogonal transformations within the information transmission context.
    \item We present several theoretical explanations for BOFT's capacity to dramatically cut down the number of learnable parameters, while retaining expressive power and stability during training. Leveraging its matrix factorization properties, BOFT also enables novel weight interpolation techniques (refer to Figure~\ref{fig:boft_interpolation}).
    \item To our knowledge, this is the initial work to extend orthogonal finetuning beyond its earlier application in controllable text-to-image synthesis~\cite{qiu2023controlling}, demonstrating BOFT's promise as a versatile method for model adaptation. We explore BOFT’s applicability to diverse adaptation tasks in both computer vision and natural language processing, consistently achieving state-of-the-art performance, which highlights its robust parameter efficiency and enhanced generalization.
\end{itemize}

\section{Related Work}

\textbf{Parameter-efficient finetuning (PEFT)}. The rapid scaling and enhancement of foundation models (\emph{e.g.}, \cite{brown2020language,radford2021learning,rombach2022high,kirillov2023segment}) has stimulated significant research into parameter-efficient strategies for adapting these models to specialized tasks~\cite{treviso2023efficient,ding2023parameter,lialin2023scaling}. Among a diverse array of PEFT techniques~\cite{houlsby2019parameter,aghajanyan2020intrinsic,hulora2022,edalati2022krona,wang2022adamix,gheini2021cross,zaken2022bitfit,guo2020parameter,sung2021training,ansell2022composable,lester2021power,li2021prefix,vu2022spot,he2021towards,mao2021unipelt,karimi2021compacter,liu2022few,sung2022lst,chen2022parameter,jia2022visual,chen2022adaptformer,zhang2022neural,jie2023fact,lian2022scaling,luo2023towards}, approaches based on reparameterization~\cite{aghajanyan2020intrinsic,hulora2022,edalati2022krona} are particularly pertinent to this study. LoRA~\cite{hulora2022}, for instance, finetunes a pretrained model by augmenting its weight matrices with products of two low-rank matrices, demonstrating strong outcomes in NLP settings. While LoRA fixes the rank of these auxiliary matrices uniformly across the architecture, subsequent approaches~\cite{zhang2022adaptive,valipour2022dylora,zhang2023increlora} permit the rank to vary between layers, thereby achieving a more adaptive parameter allocation. The straightforwardness of low-rank weight reparameterization has spurred broad adoption~\cite{dettmers2023qlora,zi2023delta,chavan2023one}. Building on insights regarding hyperspherical energy and generalization~\cite{liu2018learning,liu2021orthogonal}, \cite{qiu2023controlling} introduced orthogonal finetuning (OFT) as a promising alternative, particularly for text-to-image diffusion models: OFT employs orthogonal matrices to transform activations within individual layers, leading to improved generalization and enhanced training stability over LoRA. Nonetheless, OFT typically involves more trainable parameters compared to LoRA, motivating efforts to further reduce its parameter count. Furthermore, the generalizability of OFT to tasks beyond text-to-image model adaptation~\cite{qiu2023controlling} remains an open question. BOFT addresses OFT's parameter efficiency limitations by leveraging butterfly factorization, which, in turn, enables the broader application of orthogonal finetuning to a range of adaptation challenges.

\textbf{Butterfly structures}. The radix-2 Cooley-Tukey method decomposes an $N$-point discrete Fourier transform into two smaller transforms of size $\frac{N}{2}$, naturally leading to a recursive pattern known as the butterfly structure; mathematically, this can be represented as a product of several sparse matrices, collectively referred to as a butterfly matrix. Butterfly matrices have proven useful for representing orthogonal matrices—circumventing the need for pivoting during Gaussian elimination and facilitating efficient computations~\cite{parker1995random}, improving the stability of recurrent neural network training~\cite{jing2017tunable}, and supporting kernel approximation~\cite{munkhoeva2018quadrature}. Further, works such as~\cite{dao2019learning,dao2019kaleidoscope} have focused on learning efficient linear transformations via butterfly parameterizations, while others~\cite{chen2022pixelated,dao2022monarch} have applied butterfly matrices for computationally efficient neural network training. Additionally, butterfly structures are utilized in fast algorithms for matrix-vector multiplication~\cite{michielssen1996multilevel,o2010algorithm}, approximating data-sparse matrices~\cite{li2015butterfly}, and optimizing data transmission across networks~\cite{klasing1994broadcasting,li2003linear,rajasingh2016transmission}. Departing from prior usages, our work exploits butterfly structures specifically to improve the parameter efficiency of OFT.

\section{An Information Transmission View on Orthogonal Finetuning}

We begin by reviewing key aspects of OFT. In order to adapt a pretrained weight matrix, OFT reformulates the updated weights as the product of a trainable orthogonal matrix and the original, frozen pretrained matrix. This stands in contrast to LoRA, which introduces updates through an additive, low-rank parameterization; in OFT, the weight update instead takes the form of a multiplicative orthogonal matrix. To remain parameter-efficient, LoRA employs a low-rank decomposition, whereas the original OFT opts for a block-diagonal orthogonal parameterization~\cite{qiu2023controlling}, with parameter efficiency increasing as the diagonal blocks become smaller. Figure~\ref{fig:comp_lora} visually contrasts these methods. The underlying rationale for using orthogonal transformations in weight finetuning is to maintain the pairwise angular relationships among neurons~\cite{liu2018learning,liu2021orthogonal,qiu2023controlling}, which aids in preserving the semantic information obtained during pretraining. 

Specifically, given a pretrained linear transformation $\bm{W}^0\in\mathbb{R}^{d\times n}$, OFT learns an orthogonal matrix $\bm{R}\in\mathbb{R}^{d\times d}$ and alters the forward computation from $\bm{z}=(\bm{W}^0)^\top\bm{x}$ to $\bm{z}=(\bm{R}\bm{W}^0)^\top\bm{x}$, with $\bm{x}\in\mathbb{R}^{d}$ as the input and $\bm{z}\in\mathbb{R}^{n}$ as the output. The orthogonal matrix $\bm{R}$ is initialized as the identity matrix to ensure zero initialization. To preserve the orthogonality of $\bm{R}$ throughout training, the Cayley parameterization is employed, following \cite{qiu2023controlling,liu2021orthogonal}: $\bm{R}=(\bm{I}+\bm{Q})(\bm{I}-\bm{Q})^{-1}$, where $\bm{Q}$ is skew-symmetric, satisfying $\bm{Q}=-\bm{Q}^\top$. To enforce parameter efficiency, OFT structures $\bm{R}$ as a block-diagonal matrix: $\text{diag}(\bm{R}_1,\bm{R}_2,\cdots,\bm{R}_r)$, with each block $\bm{R}_i\in\mathcal{R}^{b\times b}, \forall i$, a smaller orthogonal matrix, and $br = d$. However, this approach incurs a limitation: the parameter-efficiency gained by block-diagonalization entails partitioning the neuron's dimensions (i.e., columns of $\bm{W}^0$) into $r$ groups, each transformed independently by a distinct orthogonal block. While effective in practice, this grouping—based solely on index—lacks theoretical justification for the structure of neuron dimensions, highlighting the appeal of a dense orthogonal matrix. This naturally leads to the following question: \emph{Is it possible to construct a dense orthogonal matrix that maintains parameter efficiency?}

In response to this issue, we suggest constructing a dense orthogonal matrix by multiplying several sparse orthogonal matrices together. To facilitate a comprehensive and accessible analysis of the sparsity structure in orthogonal matrix factorizations, we recast the task of generating a dense orthogonal matrix within the OFT framework as an information transmission process. Concretely, the construction of a dense matrix $\bm{R}\in\mathbb{R}^{d\times d}$ through the multiplication of $m$ square matrices, expressed as $\thickmuskip=2mu \medmuskip=2mu \bm{R}=\bm{B}_m\bm{B}_{m-1}\cdots\bm{B}_1$, is analogous to the propagation of information across a $d\times (m+1)$ grid of nodes, as depicted in Figure~\ref{fig:info_trans}. This information transmission perspective is inspired by the recognition that a $d$-dimensional dense square matrix encodes complete connectivity between one set of $d$ nodes and another. For instance, if $R_{ij}$ is zero in $\bm{R}$, there is no pathway for information from node $j$ to reach node $i$; if $R_{ij}$ is nonzero, such transmission is possible. Thus, expressing $\bm{R}$ as the product $\bm{B}_m\bm{B}_{m-1}\cdots\bm{B}_1$ can be interpreted as orchestrating a sequence of information exchanges governed by the individual graphs defined by each $\bm{B}_i$. Information first passes according to $\bm{B}_1$, continuing through to $\bm{B}_n$. For illustration, Figure~\ref{fig:info_trans} presents the case $\bm{R}=\bm{B}_5\bm{B}_4\bm{B}_3\bm{B}_2\bm{B}_1$, with visualizations of both their sparsity arrangements and corresponding induced graph. This graph reflects the detailed connections yielded by fully unfolding the matrix multiplications. Within this representation, each $\bm{B}_i$ is considered as defining links from nodes at layer $i$ to nodes at layer $i+1$: more precisely, the $(j_1,j_2)$ entry of $\bm{B}_i$ indicates the presence of a directed connection from the $j_2$-th node in layer $i$ to the $j_1$-th node in layer $i+1$, where a zero signifies absence of such a link. For the product $\bm{B}_5\bm{B}_4\bm{B}_3\bm{B}_2\bm{B}_1$ to achieve a fully dense form, all nodes at the initial level must have possible information paths to every node at the sixth level. If one considers only $\bm{R}=\bm{B}_4\bm{B}_3\bm{B}_2\bm{B}_1$, involving transmissions from the source nodes at the first level to recipient nodes at the fifth level, it becomes evident that certain information flows may be blocked: for example, information from node $1$ does not reach node $3$, resulting in a zero at the $(3,1)$ position in the corresponding sparsity pattern. This relationship between connectivity in the induced graph and sparsity pattern persists if we further restrict ourselves to $\bm{R}=\bm{B}_3\bm{B}_2\bm{B}_1$ or $\bm{R}=\bm{B}_2\bm{B}_1$. In the general case, ensuring that $\bm{R}\in\mathbb{R}^{d\times d}$ is dense necessitates that the collection of factors $\bm{B}_m,\cdots,\bm{B}_1$ together form a pattern of directed edges on a $d\times (m+1)$ grid, where each directed edge connects nodes between adjacent layers (columns); this construction guarantees that each node in the first layer is able to transmit information to every node at the $(m+1)$-th layer.

Adopting the information transmission perspective sheds light on the origins of sparsity within the factorization matrices used in OFT. As depicted in Figure~\ref{fig:blk_diag}, the original OFT implements a block-diagonal configuration for $\bm{R}$. While this structure decreases the parameter count, it inherently prevents $\bm{R}$ from being a dense matrix. The objective, then, is to synthesize a dense orthogonal matrix using $m$ sparse orthogonal matrices, minimizing the count of necessary trainable parameters. Within the information transmission framework, two principal criteria arise: \emph{(i)} \textbf{dense connectivity}, whereby every initial-layer node retains at least one connection path to all final-layer nodes, and \emph{(ii)} \textbf{minimum free edges}, demanding the fewest possible edges within the orthogonality constraints. It is important to recognize that orthogonality imposes nuanced restrictions on connections between successive layers. Specifically, ensuring that each $\bm{B}_i$ is full-rank—a requirement for orthogonality—necessitates $d$ edges, establishing a bijection between nodes in the $i$-th and $(i=1)$-th layers; thus, there must be no fewer than $d$ edges between any two adjacent levels (\emph{e.g.}, 4 in the scenario presented in Figure~\ref{fig:info_trans}). These mandatory $d$ edges, integral to maintaining orthogonality, should not be included in the count of trainable edges, as they are fixed (\emph{e.g.}, for a $d\times d$ orthogonal matrix where $d$ elements are nonzero, these are restricted to $\pm 1$ values). Because orthogonal matrices can be represented with fewer parameters compared to general dense matrices, this orthogonality requirement introduces further dependencies among the possible edge connections. For instance, in a $2\times 2$ orthogonal matrix, setting the $(1,1)$ entry to zero imposes that the $(2,2)$ entry must also be zero (\emph{i.e.}, the omission of one edge can necessitate the removal of another). Consequently, orthogonality may alter the permissible edge sets by forcing the addition or removal of connections. By analyzing the pattern of nonzero entries within the composed orthogonal matrix, the information transmission approach proves particularly insightful for OFT, as the focus is exclusively on the distribution of trainable nonzero elements in $\bm{R}$, rendering their actual values irrelevant.

A straightforward approach to creating a fully connected network between two layers requires $\mathcal{O}(d^2)$ connections (specifically, one dense orthogonal matrix), amounting to $d^2-d$ connections in total (e.g., $12$ connections when $d=4$). As depicted in Figure~\ref{fig:info_trans}, a possible matrix factorization example utilizes $10$ connections, which is fewer than needed for a single dense orthogonal matrix. Adopting this graphical framework allows for an investigation of OFT's parameter efficiency from a structural standpoint, making it easy to design valid matrix factorizations. This design draws on an interesting structure from the Cooley-Tukey algorithm, termed butterfly graphs~\cite{cooley1965algorithm}, which connect $d$ inputs and $d$ outputs with only $\mathcal{O}(d\log d)$ links. For instance, while the topology shown in Figure~\ref{fig:info_trans} ensures dense connectivity with $10$ links, a butterfly architecture achieves the same with just $8$. In the following section, we explain how the butterfly configuration can enhance parameter efficiency.

\section{Orthogonal Parameterization by Butterfly Factorization}

The butterfly configuration originates from the Cooley-Tukey algorithm for computing the fast Fourier transform. Within the Fourier transform, a localized alteration in the frequency spectrum can produce global modifications in the spatial domain—a process reminiscent of our information transfer challenge, where every node in the first layer must be able to reach all nodes in the final layer. This butterfly pattern is also favored as a network topology in computer systems~\cite{solihin2015fundamentals,li2011linear} due to its communication efficiency. For $k \geq 2$, where $k$ is a power of $2$, we define the \emph{butterfly factor} $\bm{B}^F(k)$ as follows:

\begin{equation}\label{eq:but_factor}
\begin{aligned}
    \bm{B}^F(k)=\begin{bmatrix}
    \text{diag}(\bm{d}_1) & \text{diag}(\bm{d}_2) \\
    \text{diag}(\bm{d}_3) & \text{diag}(\bm{d}_4)
    \end{bmatrix}\in\mathbb{R}^{k\times k},
\end{aligned}
\end{equation}

where for each $i$, the vectors $\bm{d}_i\in\mathbb{R}^{\frac{k}{2}}$. When $d=2^N$, we further specify the $d$-dimensional \emph{butterfly matrix} $\bm{B}(d)\in\mathbb{R}^{d\times d}$ through the following recursive construction:

\begin{equation}
\begin{aligned}
    \!\!\bm{B}(d)=\tilde{\bm{B}}(d,d)\!\cdot\!\begin{bmatrix}
    \bm{B}_1(\frac{d}{2})\!\!\!\!\! & \bm{0} \\
    \bm{0}\!\!\!\!\! &\bm{B}_2(\frac{d}{2})
    \end{bmatrix}= \tilde{\bm{B}}(d,d)\tilde{\bm{B}}(d,\frac{d}{2})\cdots\tilde{\bm{B}}(d,2),
\end{aligned}
\end{equation}

In this context, $\bm{B}_1(\frac{d}{2})$ and $\bm{B}_2(\frac{d}{2})$ denote two butterfly matrices, each of dimension $\frac{d}{2}$. The \emph{butterfly component} is subsequently defined as $\thickmuskip=2mu \medmuskip=2mu \tilde{\bm{B}}(d,k)=\text{diag}(\bm{B}^F_1(k),\cdots,\bm{B}^F_{d/k}(k))$, which constructs a block-diagonal matrix of dimensions $d\times d$ with block size $k$, where each $\bm{B}^F_i(k)$ is a butterfly factor as described in Equation~\ref{eq:but_factor}. This framework allows us to parameterize orthogonal matrices using butterfly matrices, provided that all multiplicative components $\tilde{\bm{B}}(d,k)$ in the overall butterfly matrix $\bm{B}(d)$ are constrained to be orthogonal. 

To illustrate, consider the specific block-diagonal matrix $\tilde{\bm{B}}(d,2)$ with block size $2$. Orthogonality of $\tilde{\bm{B}}(d,2)$ can be straightforwardly enforced by parameterizing each $2\times 2$ block via the Cayley transform, or by employing $2$-dimensional rotations, consistent with approaches from \cite{liu2021orthogonal,qiu2023controlling}. The sparsity pattern found in any butterfly component is essentially a permutation of the sparsity present in $\tilde{\bm{B}}(d,2)$, enabling the straightforward orthogonal parameterization of all butterfly components. This construction leads to an efficient representation of orthogonal matrices through a collection of $2\times 2$ orthogonal submatrices.

Extending this approach as in \cite{chen2022pixelated}, we introduce the block butterfly component $\tilde{\bm{B}}^b(d,k)$, where every entry $\bm{d}_i$ is substituted with a $b\times b$ matrix. To ensure orthogonality for the block butterfly component $\tilde{\bm{B}}^b(d,2)$, we parameterize each $2b\times 2b$ block to be orthogonal. The non-zero structure of other components $\tilde{\bm{B}}^b(d,k)$ for $k > 2$ arises via block-wise permutations of the non-zero layout in $\tilde{\bm{B}}^b(d,2)$, thereby allowing for similar orthogonal parameterizations. Integrating these elements, the forward pass in BOFT can be expressed as

\begin{equation}
    \begin{aligned}\nonumber
    \bm{z}=\big(\bm{R}(m,b)\cdot\bm{W}^0\big)^\top\bm{x},~~~~\text{s.t.}~\bigg{\{}\bm{R}(m,b)=\prod_{i=1}^m \tilde{\bm{B}}^b_{(d,i)}~~\&~~\underbrace{\big(\tilde{\bm{B}}^b_{(d,j)}\big)^\top\tilde{\bm{B}}^b_{(d,j)}=\tilde{\bm{B}}^b_{(d,j)}\big(\tilde{\bm{B}}^b_{(d,j)}\big)^\top\!=\bm{I}_d}_{\forall j\in [1, m]}\bigg{\}},
    \end{aligned}
\end{equation}

Here, for notational simplicity, $\tilde{\bm{B}}^b(d,2^{m - i+1})$ is abbreviated as $\tilde{\bm{B}}^b_{(d,i)}$, and $\bm{I}_d$ represents the $d$-dimensional identity matrix. The orthogonal matrix $\bm{R}(m,b)\in\mathbb{R}^{d\times d}$ consists of a sequence of $m$ orthogonal butterfly modules. Throughout, we refer to BOFT equipped with $\bm{R}(m,\frac{b}{2})$ as BOFT($m$,$b$), with the constraint that $b \geq 2$. In the special case where $m=1$, BOFT($1$,$b$) collapses to the block-diagonal OFT~\cite{qiu2023controlling} having block size $b$. Likewise, when $b=d$, BOFT($1$,$d$) recovers the standard OFT~\cite{qiu2023controlling} characterized by an unconstrained, dense orthogonal matrix. Using the block butterfly matrix $\bm{B}^{\frac{b}{2}}(d)$ as $\bm{R}$, specifically in BOFT($\log \frac{2d}{b}$, $b$), produces a dense orthogonal $\bm{R}$. More generally, the trainable parameter count for BOFT($m$,$b$) amounts to $\frac{1}{2}(b-1)dm$ in order to finetune a linear transformation of dimensions $d\times n$. With a butterfly matrix instantiation—i.e., setting $m=\log d, b=2$—the number of parameters scales as $\mathcal{O}(d\log d)$. By contrast, classical OFT using a full orthogonal matrix scales with $\mathcal{O}(d^2)$, while block-diagonal OFT with $r$ blocks requires $\mathcal{O}(bd)$. Thus, traditional OFT requires $b=d$ to yield a fully dense orthogonal matrix, whereas BOFT attains this property for arbitrary $b$.

\textbf{Identity initialization for BOFT}. Finetuning is typically initialized to match the original pretrained model parameters, limiting deviation during adaptation. For instance, LoRA achieves this by initializing low-rank updates to zero. In BOFT, initialization proceeds by setting all butterfly modules as identity matrices (which corresponds to zero-initialization of the skew-symmetric matrix in the Cayley parameterization).

\textbf{Multiplicative Dropout for BOFT}. LoRA~\cite{hulora2022} integrates Dropout for its low-rank updates, which effectively curbs overfitting. Standard Dropout~\cite{srivastava2014dropout} is compatible with LoRA due to its additive update structure, but is unsuitable for BOFT’s multiplicative framework. To resolve this, we design a multiplicative Dropout tailored to BOFT. As $\bm{R}(m,b)$ is made of $m$ orthogonal butterfly layers—each of which can be reorganized as $2b\times 2b$ block-diagonal orthogonal matrices—multiplicative Dropout is realized by randomly selecting a proportion $p_1$ of butterfly layers and a proportion $p_2$ of their diagonal sub-blocks, subsequently substituting those blocks by identity matrices.

\section{Intriguing Insights and Discussions}

\textbf{Expressivity of BOFT}. While the butterfly structure combined with permutations has been shown to exactly capture a variety of established fast linear transforms~\cite{dao2019learning,dao2019kaleidoscope} (such as the fast Fourier transform and Hadamard transform), the extent to which our orthogonal butterfly matrix can approximate arbitrary orthogonal matrices is not yet clear. To investigate this, we performed an empirical study by approximating a randomly generated dense orthogonal matrix of size $512\times 512$~\cite{anderson1987generation}. The results, which are averaged over 10 different random seeds, are depicted in Figure~\ref{fig:para_eff}. The vertical axis measures the approximation error, while the horizontal axis indicates the count of effective trainable parameters. Each curve of identical color corresponds to a specific block size for BOFT; the point at the far left on each curve indicates the performance of BOFT($1$,$b$), representing the original block-diagonal OFT for block size $b$. Empirically, BOFT generally achieves superior parameter efficiency relative to OFT. For instance, BOFT($9$,$2$) demonstrates greater expressivity than BOFT($1$,$16$) while utilizing substantially fewer parameters. Configurations with smaller $b$ and higher $m$ tend to be even more parameter-efficient; for example, BOFT($6$,$4$) attains similar approximation error to BOFT($2$,$16$) but with markedly fewer parameters. In summary, the butterfly matrix forms a more highly structured subset within the orthogonal group than the block-diagonal variant, making BOFT provably more expressive than OFT at the same block size.

\begin{theorem}[Expressivity of BOFT]\label{thm:exp_boft}
    BOFT surpasses OFT in expressivity when using the same block size. The full class of orthogonal matrices of dimension $d$ can be captured via products of butterfly matrices as $\bm{B}_{d-1,1}(d)\bm{B}^\top_{d-1,2}(d)\cdots\bm{B}_{1,1}(d)\bm{B}^\top_{1,2}(d)$, where every $\bm{B}_{i,j}(d),\forall i,\forall j$ denotes a butterfly matrix.
\end{theorem}

Building on Theorem~\ref{thm:exp_boft}, one arrives at a natural generalization for BOFT: the final orthogonal transformation can be extended as $\thickmuskip=2mu \medmuskip=2mu \bm{R}^G(m_1,b_1,m_2,b_2,l)=\bm{R}_{l,1}(m_1,b_1)\bm{R}_{l,2}^T(m_2,b_2)\cdots\bm{R}_{1,1}(m_1,b_1)\bm{R}_{1,2}^T(m_2,b_2)$, where $\bm{R}_{i,j}^T(m,b)$ represents the orthogonal matrix specific to the BOFT scheme. With settings $m_1=m_2=\log d$, $b_1=b_2=2$, and $l=d-1$, the generalized composition $\bm{R}^G(m_1,b_1,m_2,b_2,l)$ is expressive enough to represent any orthogonal matrix (i.e., the full orthogonal group), a structure often referred to as the kaleidoscope hierarchy~\cite{dao2019kaleidoscope}. Nonetheless, it should be noted that enhanced expressivity does not guarantee superior fine-tuning performance; even fully expressive parameterizations may not achieve optimal results after fine-tuning. This highlights the importance of balancing expressivity against regularization, which underlies the transferability of fine-tuned models. BOFT enhances the finetuning parameter landscape with built-in structural inductive biases, thereby facilitating an improved equilibrium between expressivity and regularization.

\textbf{Spectral properties}. Orthogonal finetuning consistently achieves superior spectral characteristics compared to LoRA, primarily because it exactly maintains the spectral norm of the original weight matrix $\bm{W}^0$. This invariance is clear via the singular value decomposition representation: $\bm{W}^0=\bm{U}\bm{\Sigma}\bm{V}^\top$, where $\bm{U}$ and $\bm{V}$ are orthogonal matrices, and $\bm{\Sigma}$ contains the singular values along the diagonal. Both OFT and BOFT operate by applying an orthogonal transformation $\bm{R}$ from the left, producing adapted weights $\bm{R}\bm{U}\bm{\Sigma}\bm{V}^\top$. This transformation leaves the dominant singular value (that is, the spectral norm of $\bm{W}^0$) unchanged. Preserving this property has been empirically linked to improved stability during training and enhanced generalization~\cite{miyato2018spectral,yoshida2017spectral}. Additional mathematical insights are provided in Appendix~\ref{app:sn_property}.

\textbf{Orthogonal finetuning as learning bilinear similarity}. The BOFT approach can be reformulated in terms of learning a bilinear similarity function, expressed as $\bm{w}^0_i\bm{R}\bm{x}$, with $\bm{w}^0_i$ representing the $i$-th neuron (specifically, a column in $\bm{W}^0$). In this view, BOFT is interpreted as optimizing a bilinear similarity matrix $\bm{R}$, subject to the rigorous constraint of orthogonality, thus closely relating to methodologies in distance metric learning~\cite{xing2002distance} and bilinear forms~\cite{roman2005advanced}.

\textbf{Inductive bias and generalization in BOFT}. The structure $\bm{R}(m,b)$ imposed in BOFT generally encapsulates a restricted subset within the orthogonal group, effectively constraining the possible solution space and imparting a specific inductive bias. We posit that this structured bias, originating from the butterfly factorization, enhances generalization capability; indeed, it mirrors the architecture found in a variety of foundational linear transformations~\cite{dao2019kaleidoscope} such as the discrete Fourier, discrete sine/cosine, and Hadamard transforms. Furthermore, the inherent sparsity arising from BOFT’s matrix factorization might contribute additional, perhaps implicit, inductive biases~\cite{gunasekar2017implicit,li2018algorithmic,liu2019neural}.

\textbf{Comparison to butterfly-based sparse training}. Several previous studies~\cite{dao2019learning,dao2019kaleidoscope,chen2022pixelated,dao2022monarch} have explored sparse training strategies using butterfly parameterization. These works often directly reparameterize neural network weight matrices into the butterfly format and proceed to train models from scratch. In contrast, \cite{dao2022monarch} investigates an alternative finetuning approach in which pretrained weights are initially projected onto modified butterfly matrices, followed by optimization of these projected elements for specific downstream applications. The finetuning method introduced by BOFT, however, is markedly distinct; it modifies weights using weight matrices shared across layers.

\input{rebuttal-results}

\section{Concluding Remarks and Limitations}

In this work, we introduce Orthogonal Butterfly, a universally applicable parameter-efficient finetuning approach tailored for foundation models, drawing upon the butterfly structure. Our central premise for enhanced parameter efficiency is to construct dense orthogonal matrices through the composition of multiple sparse orthogonal matrices. To facilitate practical identification of valid matrix factorizations, we develop a framework based on graphical information transmission. Within this conceptual foundation, the butterfly structure emerges as a powerful tool to realize our objectives of sparse orthogonal matrix factorization. We provide experimental evidence showcasing BOFT’s strong performance in finetuning diverse foundation models, including large language models, large vision models, and text-to-image generative systems. Our empirical results substantiate the overall effectiveness of BOFT as a versatile finetuning technique.

Nevertheless, there are limitations to BOFT’s design. Because the ultimate orthogonal matrix in BOFT is formed as a product of several orthogonal matrices, the process incurs somewhat greater training runtime overhead compared to OFT. Reducing this runtime overhead during training is an unresolved challenge. Fortunately, once finetuning concludes, the orthogonal matrices obtained via BOFT can be integrated directly into the pretrained model, introducing no extra inference delay. Additionally, it remains an open question whether the butterfly topology constitutes the most optimal structure for information transmission. The information transmission framework we establish opens avenues to draw from areas such as computer networking, where optimal network topologies for information transfer are well studied. We anticipate that adopting alternative network structures could further enhance the efficiency of composing dense orthogonal matrices.

\end{document}